% =========================================================
% Completeness Equivalences
% =========================================================

\begin{remark*}[Interpretation]
The least upper bound property is the order form of completeness. The
Monotone Convergence Theorem, the Nested Interval Property, the
Bolzano--Weierstrass theorem, and the Cauchy criterion are sequential and
interval forms of the same gap-free principle.
\end{remark*}

\begin{theorembox}{Theorem (Least Upper Bound Property Implies Monotone Convergence)}
\begin{theorem}[Least Upper Bound Property Implies Monotone Convergence]
\label{thm:lub-property-implies-monotone-convergence}
Let $F$ be a linearly ordered field satisfying the least upper bound property.
Let $(x_n)_{n\in\mathbb{N}}$ be a sequence in $F$.  If
\[
  (\forall n\in\mathbb{N})(x_n\leq x_{n+1})
\]
and
\[
  (\exists M\in F)(\forall n\in\mathbb{N})(x_n\leq M),
\]
then the set
\[
  X:=\{x_n:n\in\mathbb{N}\}
\]
has a supremum $s\in F$, and
\[
  (\forall\varepsilon\in F)\left[
  \varepsilon>0
  \Longrightarrow
  (\exists N\in\mathbb{N})(\forall n\in\mathbb{N})
  \bigl(n\geq N\Longrightarrow |x_n-s|<\varepsilon\bigr)
  \right].
\]
Consequently,
\[
  \lim_{n\longrightarrow\infty}x_n=s=\sup\{x_n:n\in\mathbb{N}\}.
\]
Dually, if
\[
  (\forall n\in\mathbb{N})(x_{n+1}\leq x_n)
\]
and
\[
  (\exists m\in F)(\forall n\in\mathbb{N})(m\leq x_n),
\]
then
\[
  \lim_{n\longrightarrow\infty}x_n=\inf\{x_n:n\in\mathbb{N}\}.
\]

\smallskip
\noindent
\hyperref[prf:lub-property-implies-monotone-convergence]{\textit{Go to proof.}}
\end{theorem}
\end{theorembox}

\begin{remark*}[Standard quantified statement]
\[
\operatorname{HasLeastUpperBoundProperty}(\mathbb{R})
\Longrightarrow
\forall (x_n)\subseteq\mathbb{R}\;
\bigl(
  \operatorname{Monotone}((x_n),\mathsf{OrderedSet}(\mathbb{R},\leq))
  \land
  \operatorname{Bounded}(x_n)
  \Longrightarrow
  \operatorname{ConvergesTo}(x_n)
\bigr).
\]
\end{remark*}

\begin{remark*}[Interpretation]
The least upper bound property supplies the missing limit for a bounded
monotone sequence: the supremum or infimum of its range becomes the sequence
limit.
\end{remark*}

\begin{dependencies}
\begin{itemize}
  \item \hyperref[def:real-sequence]{Real Sequence}
  \item \hyperref[def:sequence]{Sequence}
  \item \hyperref[def:convergent-sequence]{Convergent Sequence}
  \item \hyperref[def:monotone-sequence]{Monotone Sequence}
  \item \hyperref[def:bounded-sequence]{Bounded Sequence}
  \item \hyperref[def:least-upper-bound-property]{Least Upper Bound Property}
  \item \hyperref[thm:monotone-convergence-theorem]{Monotone Convergence Theorem}
\end{itemize}
\end{dependencies}

\begin{theorembox}{Theorem (Monotone Convergence Implies Least Upper Bound Property)}
\begin{theorem}[Monotone Convergence Implies Least Upper Bound Property]
\label{thm:monotone-convergence-implies-lub-property}
Let $F$ be an Archimedean linearly ordered field.  Suppose that every
increasing sequence in $F$ that is bounded above converges in $F$; explicitly,
suppose that for every sequence $(x_n)_{n\in\mathbb{N}}$ in $F$,
\[
  \left[(\forall n\in\mathbb{N})(x_n\leq x_{n+1})\right]
  \land
  \left[(\exists M\in F)(\forall n\in\mathbb{N})(x_n\leq M)\right]
\]
implies that there exists $L\in F$ such that
\[
  (\forall\varepsilon\in F)\left[
  \varepsilon>0
  \Longrightarrow
  (\exists N\in\mathbb{N})(\forall n\in\mathbb{N})
  \bigl(n\geq N\Longrightarrow |x_n-L|<\varepsilon\bigr)
  \right].
\]
Then $F$ has the least upper bound property.  That is, for every
$A\subseteq F$,
\[
  \left[
  A\neq\varnothing
  \land
  (\exists u\in F)(\forall a\in A)(a\leq u)
  \right]
\]
implies that there exists $s\in F$ satisfying
\[
  (\forall a\in A)(a\leq s)
\]
and
\[
  (\forall u\in F)\left[
  (\forall a\in A)(a\leq u)
  \Longrightarrow
  s\leq u
  \right].
\]

\smallskip
\noindent
\hyperref[prf:monotone-convergence-implies-lub-property]{\textit{Go to proof.}}
\end{theorem}
\end{theorembox}

\begin{remark*}[Standard quantified statement]
\[
\bigl[\forall (x_n)\subseteq\mathbb{R}\;
(\operatorname{Monotone}((x_n),\mathsf{OrderedSet}(\mathbb{R},\leq))\land\operatorname{Bounded}(x_n)
\Longrightarrow \operatorname{ConvergesTo}(x_n))\bigr]
\Longrightarrow
\operatorname{HasLeastUpperBoundProperty}(\mathbb{R}).
\]
\end{remark*}

\begin{remark*}[Interpretation]
Starting from a nonempty set $A\subseteq\mathbb{R}$ bounded above, choose an
element of $A$ and an upper bound. Bisect the bracket. Keep the upper endpoint
when the midpoint is still an upper bound; otherwise keep the midpoint as a
new lower endpoint. The two endpoint sequences are bounded and monotone, so
monotone convergence produces a common limiting boundary, which is the
supremum of $A$.
\end{remark*}

\begin{dependencies}
\begin{itemize}
  \item \hyperref[def:real-sequence]{Real Sequence}
  \item \hyperref[def:sequence]{Sequence}
  \item \hyperref[def:convergent-sequence]{Convergent Sequence}
  \item \hyperref[def:monotone-sequence]{Monotone Sequence}
  \item \hyperref[def:bounded-sequence]{Bounded Sequence}
  \item \hyperref[def:supremum]{Supremum}
  \item \hyperref[def:least-upper-bound-property]{Least Upper Bound Property}
  \item \hyperref[thm:monotone-convergence-theorem]{Monotone Convergence Theorem}
\end{itemize}
\end{dependencies}

\begin{theorembox}{Theorem (Least Upper Bound Property Equivalent to Monotone Convergence)}
\begin{theorem}[Least Upper Bound Property Equivalent to Monotone Convergence]
\label{thm:lub-property-equivalent-monotone-convergence}
Let $F$ be an Archimedean linearly ordered field.  The following statements are
equivalent:
\begin{enumerate}[label=(\roman*)]
  \item Every nonempty subset of $F$ that is bounded above has a least upper
  bound in $F$.
  \item Every increasing sequence in $F$ that is bounded above converges in
  $F$.
  \item Every decreasing sequence in $F$ that is bounded below converges in
  $F$.
\end{enumerate}

\smallskip
\noindent
\hyperref[prf:lub-property-equivalent-monotone-convergence]{\textit{Go to proof.}}
\end{theorem}
\end{theorembox}

\begin{remark*}[Standard quantified statement]
\[
\operatorname{HasLeastUpperBoundProperty}(\mathbb{R})
\Longleftrightarrow
\forall (x_n)\subseteq\mathbb{R}\;
\bigl(\operatorname{Monotone}((x_n),\mathsf{OrderedSet}(\mathbb{R},\leq))\land\operatorname{Bounded}(x_n)
\Longrightarrow \operatorname{ConvergesTo}(x_n)\bigr).
\]
\end{remark*}

\begin{remark*}[Interpretation]
The equivalence identifies order completeness and monotone sequential
completeness as the same principle expressed in different languages.
\end{remark*}

\begin{dependencies}
\begin{itemize}
  \item \hyperref[def:real-sequence]{Real Sequence}
  \item \hyperref[def:sequence]{Sequence}
  \item \hyperref[def:convergent-sequence]{Convergent Sequence}
  \item \hyperref[def:monotone-sequence]{Monotone Sequence}
  \item \hyperref[def:bounded-sequence]{Bounded Sequence}
  \item \hyperref[def:least-upper-bound-property]{Least Upper Bound Property}
  \item \hyperref[thm:lub-property-implies-monotone-convergence]{Least Upper Bound Property Implies Monotone Convergence}
  \item \hyperref[thm:monotone-convergence-implies-lub-property]{Monotone Convergence Implies Least Upper Bound Property}
\end{itemize}
\end{dependencies}

\begin{theorembox}{Theorem (Least Upper Bound Property Equivalent to Nested Intervals)}
\begin{theorem}[Least Upper Bound Property Equivalent to Nested Intervals]
\label{thm:lub-property-equivalent-nested-interval-property}
Let $F$ be an Archimedean linearly ordered field.  The following statements are
equivalent:
\begin{enumerate}[label=(\roman*)]
  \item Every nonempty subset of $F$ that is bounded above has a least upper
  bound in $F$.
  \item For every pair of sequences $(a_n)$ and $(b_n)$ in $F$ satisfying
  \[
    (\forall n\in\mathbb{N})(a_n\leq b_n)
  \]
  and
  \[
    (\forall n\in\mathbb{N})([a_{n+1},b_{n+1}]\subseteq [a_n,b_n]),
  \]
  there exists $x\in F$ such that
  \[
    (\forall n\in\mathbb{N})(a_n\leq x\leq b_n).
  \]
  Equivalently,
  \[
    \bigcap_{n=1}^{\infty}[a_n,b_n]\neq\varnothing.
  \]
\end{enumerate}

\smallskip
\noindent
\hyperref[prf:lub-property-equivalent-nested-interval-property]{\textit{Go to proof.}}
\end{theorem}
\end{theorembox}

\begin{remark*}[Standard quantified statement]
\[
\operatorname{HasLeastUpperBoundProperty}(\mathbb{R})
\Longleftrightarrow
\text{every nested sequence of nonempty closed bounded intervals has nonempty intersection}.
\]
\end{remark*}

\begin{remark*}[Interpretation]
Nested intervals express completeness geometrically: repeatedly shrinking
closed bounded intervals cannot shrink around a missing point in the real
line.
\end{remark*}

\begin{dependencies}
\begin{itemize}
  \item \hyperref[def:sequence]{Sequence}
  \item \hyperref[def:closed-interval]{Closed Interval}
  \item \hyperref[def:bounded]{Bounded Set}
  \item \hyperref[def:least-upper-bound-property]{Least Upper Bound Property}
  \item \hyperref[thm:nested-interval-theorem]{Nested Interval Theorem}
\end{itemize}
\end{dependencies}

\begin{theorembox}{Theorem (Standard Completeness Equivalences)}
\begin{theorem}[Standard Completeness Equivalences]
\label{thm:standard-completeness-equivalences}
Let $F$ be an Archimedean linearly ordered field.  The following statements are
equivalent:
\begin{enumerate}[label=(\roman*)]
  \item $F$ has the least upper bound property.
  \item $F$ has the greatest lower bound property.
  \item Every increasing sequence in $F$ that is bounded above converges, and
  every decreasing sequence in $F$ that is bounded below converges.
  \item Every nested sequence of nonempty closed bounded intervals in $F$ has
  nonempty intersection.
  \item Every Cauchy sequence in $F$ converges in $F$.
  \item Every bounded sequence in $F$ has a convergent subsequence.
\end{enumerate}

\smallskip
\noindent
\hyperref[prf:standard-completeness-equivalences]{\textit{Go to proof.}}
\end{theorem}
\end{theorembox}

\begin{remark*}[Standard quantified statement]
\[
\text{LUB property}
\Longleftrightarrow
\text{monotone convergence}
\Longleftrightarrow
\text{nested intervals}
\Longleftrightarrow
\text{Bolzano--Weierstrass}
\Longleftrightarrow
\text{Cauchy completeness}.
\]
\end{remark*}

\begin{remark*}[Interpretation]
The first statement is the order version. The second and fifth are sequence
versions. The third is the interval version. The fourth is the compactness
version for bounded sequences.
\end{remark*}

\begin{dependencies}
\begin{itemize}
  \item \hyperref[def:real-sequence]{Real Sequence}
  \item \hyperref[def:sequence]{Sequence}
  \item \hyperref[def:subsequence-of-sequence]{Subsequence}
  \item \hyperref[def:convergent-sequence]{Convergent Sequence}
  \item \hyperref[def:closed-interval]{Closed Interval}
  \item \hyperref[def:cauchy-sequence]{Cauchy Sequence}
  \item \hyperref[def:monotone-sequence]{Monotone Sequence}
  \item \hyperref[def:bounded-sequence]{Bounded Sequence}
  \item \hyperref[def:least-upper-bound-property]{Least Upper Bound Property}
  \item \hyperref[thm:monotone-convergence-theorem]{Monotone Convergence Theorem}
  \item \hyperref[thm:nested-interval-theorem]{Nested Interval Theorem}
  \item \hyperref[thm:bolzano-weierstrass-sequences]{Bolzano--Weierstrass Theorem for Sequences}
  \item \hyperref[thm:cauchy-criterion-real-sequences]{Cauchy Criterion for Real Sequences}
\end{itemize}
\end{dependencies}
